\section{Closed-half-line smoothing}\label{v9:prawitz:appendix}
The finite-index reduction in Appendix~\ref{r4bridge:section} needs a
majorant for the closed event, including equality at an atom. We give its
construction and normalization here. This appendix is independent of the
Fourier remainder used in the fixed-target coarse entry.

\subsection{Exact half-line smoothing}
\label{ms:fourier:kernel}

We construct a majorant of the closed half-line and justify its
expectation under an integrable random variable. Its value at the boundary
is part of the proof, because the finite-index laws can have atoms.

Write $\phi$ and $\Phi$ for the standard normal density and distribution
function. All terminating decimals are exact rational numbers. The
imaginary unit is $i$. Symbols introduced within a subsection, including
$D$, $R$, $L$, $s$ and $K$, are local to that subsection unless explicitly
carried forward. In particular the moment defect $D$, an auxiliary distance
bound, and the Bernoulli inward displacement $d$ are different quantities.

The one-sided kernel is
\[
 K(t)=\frac{1-t}{2}+\frac{i}{2}
       \left((1-t)\cot(\pi t)+\frac1\pi\right),\qquad 0<t<1.
\]
For $\kappa\in\{K,\overline K\}$, $0<b\le1$ and
$g(t)=g(0)+O(t)$ with real $g(0)$, the convention throughout is
\begin{equation}\label{ms:fourier:kernel-convention}
 \mathfrak I_{\kappa,b}(g)
 =2\lim_{\epsilon\downarrow0}\operatorname{Re}
       \int_\epsilon^b\kappa(t)g(t)\,dt.
\end{equation}
Only the real-part improper limit is intended; the unmodified complex
integral need not converge. The proof below establishes directly that
\[
 M(x)=\tfrac12+\mathfrak I_{K,1}(e^{2\pi itx})
       \ge\mathbf1_{(-\infty,0]}(x),\qquad
 1-M(-x)\le\mathbf1_{(-\infty,0)}(x),
\]
including atoms at the boundary, and justifies the expectation identity
for an integrable random variable.

\subsubsection{A real-variable Beurling construction}

Set

\[
 W(x)=\left(\frac{\sin\pi x}{\pi x}\right)^2\quad(x\ne0),
 \qquad W(0)=1,
\]

and define the continuous odd function

\begin{equation}
 H_B(x)=\sum_{m=1}^{\infty}\{W(x-m)-W(x+m)\}+2xW(x).
 \label{ms:fourier:kernel-1}
\end{equation}

The series converges absolutely and locally uniformly on the real line:
for \(|x|\le R\) and \(m>2R+1\), each shifted term is bounded by a
constant times \(m^{-2}\). Thus the function and all the integer values in
what follows are unambiguous. In particular,

\begin{equation}
 H_B(0)=0,\qquad H_B(j)=\operatorname{sgn}(j)
 \quad(j\in\mathbb Z\setminus\{0\}).                 \label{ms:fourier:kernel-2}
\end{equation}

Here \(\operatorname{sgn}(0)=0\). The corresponding classical Beurling
majorant is \(B=H_B+W\). For noninteger \(x\), its series is

\begin{equation}
 B(x)=\left(\frac{\sin\pi x}{\pi}\right)^2
 \left\{\frac2x+\sum_{m=0}^{\infty}\frac1{(x-m)^2}
                    -\sum_{m=1}^{\infty}\frac1{(x+m)^2}\right\}.
 \label{ms:fourier:kernel-3}
\end{equation}

We prove the needed inequality directly, without importing an extremality
or Fourier-transform theorem about \(B\).

The elementary partial-fraction identity

\begin{equation}
 \sum_{m\in\mathbb Z}\frac1{(y-m)^2}
       =\frac{\pi^2}{\sin^2\pi y},\qquad y\notin\mathbb Z,
 \label{ms:fourier:kernel-4}
\end{equation}

has the following short proof. Integrate
\(\pi\cot(\pi z)/(z-y)^2\) counterclockwise over the square with vertices
\(\pm(N+\tfrac12)\pm i(N+\tfrac12)\). The cotangent is uniformly bounded
on these edges: on the vertical edges it is a multiple of a hyperbolic
tangent, and on the horizontal edges its modulus is bounded by
\(\coth(\pi(N+\tfrac12))\). The denominator is of order \(N^2\) and the
perimeter of order \(N\), so the integral tends to zero. Its residues are
\((m-y)^{-2}\) at \(m=-N,\ldots,N\), and
\(-\pi^2\csc^2(\pi y)\) at \(y\). The residue theorem and absolute
convergence give \eqref{ms:fourier:kernel-4}.

For \(y>0\), \(y\notin\mathbb Z\), put

\[
 A(y)=\left(\frac{\sin\pi y}{\pi}\right)^2,
 \qquad S(y)=\sum_{m=1}^{\infty}(y+m)^{-2}.
\]

Subtracting \eqref{ms:fourier:kernel-4} from \eqref{ms:fourier:kernel-1} gives exactly

\begin{equation}
 H_B(y)-1=A(y)\left\{\frac2y-\frac1{y^2}-2S(y)\right\}.
 \label{ms:fourier:kernel-5}
\end{equation}

The decreasing positive function \(s\mapsto(y+s)^{-2}\) satisfies

\begin{equation}
 S(y)\le\int_0^\infty(y+s)^{-2}\,ds=\frac1y
       \le\frac1{y^2}+S(y).
 \label{ms:fourier:kernel-6}
\end{equation}

Consequently the expression in braces in \eqref{ms:fourier:kernel-5} lies between
\(-y^{-2}\) and \(y^{-2}\). Hence

\[
 |H_B(y)-1|\le W(y)\quad(y>0).
\]

Positive integer points follow directly from \eqref{ms:fourier:kernel-2}. Oddness of \(H_B\) and
evenness of \(W\), together with their values at zero, now prove

\begin{equation}
 |H_B(x)-\operatorname{sgn}(x)|\le W(x)
 \qquad(x\in\mathbb R).                              \label{ms:fourier:kernel-7}
\end{equation}

This is the required sign-function majorization, with its normalization
proved rather than assumed.

\subsubsection{Recovering the exact Fourier kernel}

Elementary integration gives, for every real \(x\),

\begin{equation}
 W(x)=2\int_0^1(1-t)\cos(2\pi tx)\,dt.              \label{ms:fourier:kernel-8}
\end{equation}

For \(0<r<1\), insert \eqref{ms:fourier:kernel-8} into the absolutely convergent Abel-weighted
series. Termwise integration is legitimate because
\(\sum_{m\ge1}r^m<\infty\). The cosine subtraction formula gives

\begin{equation}
 \begin{aligned}
 \sum_{m\ge1}r^m\{W(x-m)-W(x+m)\}
   &=4\int_0^1(1-t)\sin(2\pi tx)
               \sum_{m\ge1}r^m\sin(2\pi mt)\,dt\\
   &=4\int_0^1(1-t)\sin(2\pi tx)
        \frac{r\sin(2\pi t)}{1-2r\cos(2\pi t)+r^2}\,dt.
 \end{aligned}                                                    \label{ms:fourier:kernel-9}
\end{equation}

For \(0<t<1\),

\begin{equation}
 \left|\frac{r\sin(2\pi t)}{1-2r\cos(2\pi t)+r^2}\right|
 \le\frac12|\cot(\pi t)|,
 \label{ms:fourier:kernel-10}
\end{equation}

since the denominator is
\((1-r)^2+4r\sin^2(\pi t)\). The right side of \eqref{ms:fourier:kernel-9}, including its
remaining factors, is therefore dominated by

\[
 2(1-t)|\sin(2\pi tx)|\,|\cot(\pi t)|.
\]

This is integrable at both endpoints: the sine is \(O_x(t)\) at zero,
and \((1-t)\cot(\pi t)\) is bounded at one. Dominated convergence on
the integral side and absolute convergence on the series side allow
\(r\uparrow1\). The fraction in \eqref{ms:fourier:kernel-9} tends to
\(\tfrac12\cot(\pi t)\). Also, direct integration gives

\begin{equation}
 2xW(x)=\frac2\pi\int_0^1\sin(2\pi tx)\,dt,          \label{ms:fourier:kernel-11}
\end{equation}

including \(x=0\). Combining \eqref{ms:fourier:kernel-1} and \eqref{ms:fourier:kernel-9}–\eqref{ms:fourier:kernel-11},

\begin{equation}
 H_B(x)=2\int_0^1 R(t)\sin(2\pi tx)\,dt,
 \qquad R(t)=(1-t)\cot(\pi t)+\frac1\pi.
 \label{ms:fourier:kernel-12}
\end{equation}

Every integral in \eqref{ms:fourier:kernel-12} is a convergent real integral. No principal-value
Fourier manipulation has been used.

Now take exactly the manuscript's kernel

\begin{equation}
 K(t)=\frac{1-t}{2}+\frac i2 R(t),\qquad0<t<1.
 \label{ms:fourier:kernel-13}
\end{equation}

Equations \eqref{ms:fourier:kernel-8} and \eqref{ms:fourier:kernel-12} establish the exact identity

\begin{equation}
 \boxed{\displaystyle
 M(x):=\frac12+2\lim_{\epsilon\downarrow0}\operatorname{Re}
       \int_\epsilon^1K(t)e^{2\pi itx}\,dt
       =\frac{1-H_B(x)+W(x)}2=\frac{1+B(-x)}2.}
 \label{ms:fourier:kernel-14}
\end{equation}

In particular, the exponent has the \emph{positive} sign and the imaginary
part of \(K\) has the \emph{positive} sign. The minus sign on \(H_B\) comes
from taking the real part of \(iR(t)e^{2\pi itx}\).

\subsubsection{Closed boundary, cardinal values and the minorant}

For \(x<0\), \eqref{ms:fourier:kernel-7} gives \(M(x)\ge1\); for \(x>0\), it gives
\(M(x)\ge0\). At the boundary, \eqref{ms:fourier:kernel-2} and \(W(0)=1\) give

\begin{equation}
 M(0)=1.                                             \label{ms:fourier:kernel-15}
\end{equation}

Thus this majorizes the \emph{closed} half-line, including a point mass at
zero, rather than merely the half-height convention for a jump:

\begin{equation}
 M(x)\ge1_{(-\infty,0]}(x),\qquad
 1-M(-x)\le1_{(-\infty,0)}(x).
 \label{ms:fourier:kernel-16}
\end{equation}

The second inequality follows by subtracting the first inequality at
\(-x\) from one; its endpoint is open. Direct cardinal values are

\begin{equation}
 M(j)=1_{\{j\le0\}},\qquad1-M(-j)=1_{\{j<0\}}
 \quad(j\in\mathbb Z).                              \label{ms:fourier:kernel-17}
\end{equation}

The proof also gives the useful global error bound

\begin{equation}
 0\le M(x)-1_{(-\infty,0]}(x)\le W(x)
 \quad(x\in\mathbb R).                              \label{ms:fourier:kernel-18}
\end{equation}

At zero the left-hand error is exactly zero; it is not obtained by
assigning half of an atom to each side. In particular \(0\le M\le2\).

\subsubsection{Expectations and the compensated identity}

Suppose \(U\) is real and \(\E|U|<\infty\), and use the convention
\(\psi_U(t)=\E e^{2\pi itU}\). Since
\(2\int_0^1\operatorname{Re}K(t)\,dt=1/2\), \eqref{ms:fourier:kernel-14} can be written

\begin{equation}
 M(u)=1+2\operatorname{Re}\int_0^1
             K(t)\{e^{2\pi itu}-1\}\,dt.             \label{ms:fourier:kernel-19}
\end{equation}

The complex integral in \eqref{ms:fourier:kernel-19} is absolutely convergent. Indeed,
\(K(t)=i/(2\pi t)+1/2+O(t)\) at zero, \(K\) is bounded near one,
and

\[
 \E|e^{2\pi itU}-1|\le2\pi t\,\E|U|.
\]

Consequently
\(\int_0^1|K(t)|\E|e^{2\pi itU}-1|\,dt<\infty\), so ordinary
absolute Fubini applies to \eqref{ms:fourier:kernel-19}. Restoring the constant real part gives

\begin{equation}
 \E M(U)=\frac12+
     2\lim_{\epsilon\downarrow0}\operatorname{Re}
        \int_\epsilon^1K(t)\psi_U(t)\,dt.            \label{ms:fourier:kernel-20}
\end{equation}

Since \(\psi_{-U}=\overline{\psi_U}\), conjugation similarly yields

\begin{equation}
 \E\{1-M(-U)\}=\frac12-
     2\lim_{\epsilon\downarrow0}\operatorname{Re}
        \int_\epsilon^1\overline{K(t)}\psi_U(t)\,dt.
 \label{ms:fourier:kernel-21}
\end{equation}

These are the two expectation identities under the real-part
improper convention \eqref{ms:fourier:kernel-convention}. All expectations are finite; \eqref{ms:fourier:kernel-18} even gives
boundedness of the majorant without a moment assumption. The first
absolute moment is a sufficient hypothesis for the stated compensated
Fourier/Fubini argument. The uncompensated complex integral itself need
not converge.

\subsubsection{Four-summand smoothing normalization}

Let \(Y\) be integrable, let \(h\in\mathbb R\), and let \(T>0\). For
\(\chi_Y(s)=\E e^{isY}\), apply \eqref{ms:fourier:kernel-16} and \eqref{ms:fourier:kernel-20} to
\(U=T(Y-h)/(2\pi)\). The resulting exact normalization is

\begin{equation}
 \Pp(Y\le h)\le\frac12+
  2\lim_{\epsilon\downarrow0}\operatorname{Re}
       \int_\epsilon^1K(t)e^{-iTth}\chi_Y(Tt)\,dt.
 \label{ms:fourier:kernel-22}
\end{equation}

For the four-summand application,
\(Z=(X_1+\cdots+X_4)/2\), \(\varphi_P(s)=\E e^{isX_1}\),
and iid summands, choose \(T=6\). Independence gives
\(\chi_Z(6t)=\varphi_P(3t)^4\). Thus \eqref{ms:fourier:kernel-22} is exactly

\begin{equation}
 F_Z(x)\le\frac12+\int_0^1\operatorname{Re}
   \{\widehat K(t)e^{-6itx}\varphi_P(3t)^4\}\,dt,
 \qquad \widehat K=2K,
 \label{ms:fourier:kernel-23}
\end{equation}

the physical four-summand smoothing normalization. Writing \(D=x-Z\) and
\(S(t)=tR(t)\), its real integrand is exactly

\begin{equation}
 (1-t)\E\cos(6Dt)+S(t)\E[6D\operatorname{sinc}(6Dt)],
 \qquad \operatorname{sinc}(v)=\sin(v)/v,
 \label{ms:fourier:kernel-24}
\end{equation}

with the continuous value at zero. This verifies both the factor six
and the sign of the sine term. No non-atomicity, finite-support, density,
or continuity-at-the-threshold assumption is needed. For \eqref{ms:fourier:kernel-23}, iid
independence and integrability are sufficient; the manuscript's bounded
support assumptions imply integrability.

